\begin{wideslide}{Buffon's needle problem}
\null\vfill

  \twocolumn
  {
    {\it Suppose we have a floor made of parallel strips of wood, each the same width, and we drop a needle onto the floor. What is the probability that the needle will lie across a line between two strips?}
    
    \sep
    
    {\it\hfill Georges-Louis Leclerc,\\\hfill Comte de Buffon\\\hfill 18th century}
  }
  {
    \sep\sep
    \centering\input{figures/buffon.tex}
  }
  
  \vspace{-10pt}
  \myBoxFullWidth{Monte Carlo without computers}
  
  \twocolumn
  {
    If needle length ($l$) $<$ lines width ($t$):
    
    $$P = \frac{2l}{t\pi}$$
    
    which can be used to estimate $\pi$:
    
    $$\pi = \frac{2l}{tP}$$
  }
  {
    MC experiment was performed by Mario Lazzarini in 1901 by throwing 3408 needles:
    
    $$\pi = \frac{2l \cdot 3408}{t \cdot \#red} = \frac{355}{113} = 3.14159292$$
  }
  
\vfill\null
\end{wideslide}

\begin{wideslide}[toc = From Solitaire to MC]{From Solitaire to Monte Carlo method}
\null\vfill

    \twocolumn
    {
      \begin{itemize}
	\item Stanis{\l}aw Ulam was a Polish mathematician
	\item He invented the Monte Carlo method while playing solitaire
	\item The method was used in Los Alamos, performed by ENIAC computer
      \end{itemize}  
      
      \includegraphics[width=\columnwidth]{figures/eniac1946.eps}
    }
    {
      \includegraphics[width=\columnwidth]{figures/solitaire.eps}
      \begin{itemize}
	\item What is a probability of success in solitaire?
	\begin{itemize}
	  \item Too complex for an analytical calculations
	  \item Lets try $N = 100$ times and count wins
	  \item With $N \rightarrow \infty$ we are getting closer to correct result
	\end{itemize}
      \end{itemize}
    }
 
\vfill\null
\end{wideslide}

\begin{slide}{Newton-Pepys problem}
\null\vfill

  \twocolumn
  {
    {\it
    Which of the following three propositions has the greatest chance of success?
    
    \begin{itemize}
      \item[A] Six fair dice are tossed independently and at least one “6” appears.
      \item[B] Twelve fair dice are tossed independently and at least two “6”s appear.
      \item[C] Eighteen fair dice are tossed independently and at least three “6”s appear.
    \end{itemize}
    }
  }
  {
    \sep\sep\sep\sep
    \centering\scalebox{0.15}{\input{figures/dice.tex}}
  }

\vfill\null
\end{slide}

\begin{slide}[toc=]{Newton-Pepys problem: analytical attempt}
\null\vfill

  \begin{itemize}
    \item First, lets go back to high school and calculate this analytically
    \item Let $p = \frac{1}{6}$ be the probability of rolling $6$
    \item The probability of not rolling $6$ is $(1 - p)$
    
    \item[A] six attempts, at least one six

    $$P_A = 1 - (1 - p)^6 \approx 0.6651$$

    \vspace{-5pt}
    \item[B] twelve attempts, at least two sixes
    
    $$P_B = 1 - (1 - p)^{12} - {12 \choose 1}p(1 - p)^{11} \approx 0.6187$$

    \vspace{-5pt}
    \item[C] eighteen attempts, at least three sixes
    
    $$P_C = 1 - (1 - p)^{18} - {18 \choose 1}p(1 - p)^{17} - {18 \choose 2}p^2(1 - p)^{16} \approx 0.5973$$
  
  \end{itemize}

\vfill\null
\end{slide}

\begin{wideslide}[method=direct,toc=]{Newton-Pepys problem: MC attempt}
\null\vfill

  \begin{minipage}{0.45\textwidth}
    \begin{itemize}
     \item MC attempt is just ``performing the experiment'', so we will be rolling dices
     \item Roll $6n$ times and check if number of sixes is greater or equal $n$
     \item Repeat $N$ times and your probability is given by:
     
     $$P = \frac{\mbox{number of successes}}{N}$$
    \end{itemize}
  \end{minipage}\hspace{0.05\textwidth}\begin{minipage}{0.5\textwidth}
{\small
\begin{verbatim}
def throw (nSixes):
  n = 0
  for _ in range (6 * nSixes):
    if random.randint (1, 6) == 6: n += 1
  return n >= nSixes
  
def MC (nSixes, nAttempts):
  n = 0
  for _ in range (nAttempts):
    n += throw (nSixes)
  return float (n) / nAttempts

if __name__ == "__main__":
  for i in range (1, 4):
    print MC (i, 1000)
\end{verbatim}
}
  \end{minipage}
  
\vfill\null
\end{wideslide}

\begin{slide}[toc=]{Newton-Pepys problem: summary}
\null\vfill

  \begin{itemize}
    \item Your MC result depends on $N$
    \item Results for $N = 100$:
    
    \begin{eqnarray*}
      P_A & = & 0.71, 0.68, 0.76, 0.65, 0.68 \hspace{20pt}{\color{pdcolor3}P_A^{true} = 0.6651} \\
      P_B & = & 0.70, 0.56, 0.60, 0.63, 0.69 \hspace{20pt}{\color{pdcolor3}P_B^{true} = 0.6187}\\
      P_C & = & 0.62, 0.62, 0.53, 0.57, 0.62 \hspace{20pt}{\color{pdcolor3}P_C^{true} = 0.5973}
    \end{eqnarray*}
    
    \item Results for $N = 10^6$:
    
    \begin{eqnarray*}
      P_A & = & 0.6655, 0.6648, 0.6653, 0.6662, 0.6653 \\
      P_B & = & 0.6188, 0.6191, 0.6191, 0.6190, 0.6182 \\
      P_C & = & 0.5975, 0.5979, 0.5972, 0.5978, 0.5973
    \end{eqnarray*}
    
    \item Your MC results also depends on the way how random numbers were generated
    
  \end{itemize}

\vfill\null
\end{slide}

%\begin{wideslide}[toc = Simple Monte Carlo, method=direct]{Simple Monte Carlo - coin flipping}
%\null\vfill
%
%  \begin{itemize}
%    \item What is the probability of getting head / tail when throwing a coin?
%  \end{itemize}
%  
%  \begin{minipage}{0.7\textwidth}
%  \begin{itemize}
%    \item One can use MC method to calculate this:
%    \begin{itemize}
%      \item get $N$ times a random number $x \in \{0,1\}$
%      \item check how many ($n$) times $x = 1$
%      \item $P = \frac{n}{N}$
%    \end{itemize}
%    \item One-line example in Python:
%  \end{itemize}
%  \end{minipage}\begin{minipage}{0.3\textwidth}
%		  \includegraphics[width=0.5\columnwidth]{figures/coin.eps}
%                \end{minipage}
%
%  {\small\color{pdcolor3}
%  \begin{verbatim}
%MC = lambda N: len(filter(lambda x: random.randint(0,1), range(N))) / float(N)     
%  \end{verbatim}
%  }
%  \begin{minipage}{0.55\textwidth}
%  \begin{itemize}
%   \item If you do not like one-liners:
%  \end{itemize}
%  \vspace{-20pt}  
%  {\small\color{pdcolor3}
%  \begin{verbatim}
%def MonteCarlo(N):
%  n = 0
%  for _ in range(N):
%    if random.randint(0,1): n += 1
%  return float(n) / N   
%  \end{verbatim}
%  }
%  \end{minipage}\begin{minipage}{0.45\textwidth}
%		  \vspace{-10pt}
%		  \begin{itemize}
%		   \item Some results for MC(100):
%		   \item[]
%		   0.51, 0.48, 0.61, 0.46, 0.57, 0.46
%		   \item Some results for MC(100000):
%		   \item[]
%		   0.50139, 0.49798, 0.49933, 0.5005
%		  \end{itemize}
%                \end{minipage}
%
%\vfill\null
%\end{wideslide}